\clearpage
\hypertarget{native-472-august}{}
\begin{center}
\rule{0.55\linewidth}{0.4pt}\\[2pt]
Astronomische Nachrichten. Vol. XXII. No. 513. pp. 141..144. 22 August 1844.\\[2pt]
\rule{0.55\linewidth}{0.4pt}
\end{center}
\begin{flushright}Göttingen, 8 Aug. 1844.\end{flushright}

Since the 23rd of last month I have indeed seen the comet briefly on several evenings, but I was only able to obtain another observation yesterday. Clouds returned as early as the third comparison, making further observations impossible. Here is what I obtained, using two stars from \textsc{Bessel}’s zone no. 364 for comparison.

\begin{center}
\setlength{\tabcolsep}{3pt}
\begin{tabular}{llll}
1844 Aug. 7. & $9^{\rm u}\ 16^{\rm m}\ 29^{\rm s}$ mean time & R.A. & $=213^\circ\ 52'\ 6''8$ \\
 & $9\ 24\ 48\mathord{\cdot}$\quad --- & Decl. & $=+24\ 50\ 21.5$ \\
\end{tabular}
\end{center}

{\rightskip=0pt plus 1em By a rough estimate, Dr. \textsc{Goldschmidt}’s parabolic elements give a right ascension $10''$ greater and a declination $16''$ smaller.\par}

\begin{flushright}Göttingen, 10 Aug. 1844.\end{flushright}

I am sending you another determination of the comet’s position, which I made yesterday evening. I used 2 stars from \textsc{Bessel}’s zone 364 for comparison, the first of which is also in zone 460.

\begin{center}
\setlength{\tabcolsep}{3pt}
\begin{tabular}{llll}
1844 Aug. 9. & $9^{\rm u}\ 36^{\rm m}\ 54^{\rm s}$ mean time & G. R. & $212^\circ\ 40'\ 23''8$ \\
 & $9\ 35\ 8$\quad --- & Decl. & $23\ 18\ 16.6$ \\
\end{tabular}
\end{center}

\textsc{Goldschmidt}’s elements give a right ascension $5''$ greater and a declination $16''$ smaller.
